\section{Absolutely continuous spectrum}%\label{secxavier}
%%%%%%%%%%%%%%%%%%%%%%%%%%%%%%%%%%%%%%%%%%%%%%

Let $S$ be a self-adjoint operator in a separable Hilbert space $H$,
and denote by $R(z)=(S-z)^{-1}$ the resolvent of $S$.
Say that a closed operator $T$ in $H$ is \emph{$S$-smooth} if, for each $x\in H$ and $\ve\ne0$,
$R(\lambda+{\rm i}\ve)x\in\dom T$
 for almost all $\lambda\in\R$ and
\begin{align*}%\label{Kato1}
 \sup_{|x|=1,\,\ve\ne0} \int_{-\infty}^\infty |TR(\lambda+{\rm i}\ve)x|^2 < \infty,
\end{align*}
see~\cite[p.~142]{ReedSimon78}.
In our context, the point is that then the image
\begin{align}\label{Kato3}
 \ran T^* \subseteq H_{\ac}(S),
\end{align}
see~\cite[Theorem XIII.23]{ReedSimon78}.
For a Borel subset $B\subseteq\R$,
say that a closed operator $T$ in $H$ is \emph{$S$-smooth on $B$} if $TE_B$ is $H$-smooth,
where $E_B$ denotes the spectral projection of $S$ associated to $B$.
Then
\begin{align*}%\label{Kato5}
 \ran E_BT^* \subseteq H_{\ac}(S),
\end{align*}
by \eqref{Kato3}.
By~\cite[Theorem XIII.30]{ReedSimon78}, $T$ is $S$-smooth on the closure $\bar B$ of $B\subseteq\R$ if
\begin{align}\label{Kato7}
 \dom T\supseteq\dom S
 \qquad\text{and}\qquad
 \sup_{\substack{|x|=1, \, \lambda\in B,\\ 0<|\ve|<1}} |\ve||TR(\lambda+{\rm i}\ve)x| < \infty.
\end{align}
Xavier~\cite[p.~581f.]{Xavier1988} shows that the latter criterion is satisfied
if $B$ is bounded and if there is a~constant $C$ such that
\begin{align}\label{Kato9}
 |Tx|^2 \le C|(S-\lambda)x|(|Sx|+|x|)
\end{align}
for all $x$ in a core of $S$ in $H$ and all $\lambda\in B$.
In fact, he shows that the term in \eqref{Kato7} is then bounded by $C(1+|\lambda|+|\ve|)$,
which explains why $B$ is assumed to be bounded; see~\cite[top of~p.~582]{Xavier1988}.

Guided by~\cite[Section 3]{Xavier1988}, we return to the geometric situation in Section~\ref{secrell}
and let $X$ be a~$C^1$ vector field and $u$ be a $C^2$ function on a Riemannian manifold $M$
such that $\supp X\cap\supp u$ is compact.
Then \eqref{rell} gives, for any smooth domain $D$ in $M$ with outer normal $\nu$ along $\partial D$,
\begin{gather*}
 \int_D\langle X,\nabla u\rangle\Delta u + \frac12\int_D|\nabla u|^2\diver X \\
 \qquad\quad{}= \int_D\langle\nabla_{\nabla u}X,\nabla u\rangle
 - \int_{\partial D}\biggl(\langle X,\nabla u\rangle\langle\nabla u,\nu\rangle
 - \frac12|\nabla u|^2\langle X,\nu\rangle\biggr).
\end{gather*}
As in~\cite[p.~583]{Xavier1988}, we use
\begin{align*}
 |\nabla u|^2 = u\Delta u - \frac12\Delta u^2
\end{align*}
and obtain
\begin{align*}
 \int_D \Delta u\biggl(\langle X,\nabla u\rangle + \frac{u}2\diver X\biggr)
 = {}&\int_D \biggl(\langle\nabla_{\nabla u}X,\nabla u\rangle + \frac14\diver X\Delta u^2\biggr) \\
 &{}- \int_{\partial D}\biggl(\langle X,\nabla u\rangle\langle\nabla u,\nu\rangle
 - \frac12|\nabla u|^2\langle X,\nu\rangle\biggr).
\end{align*}
Since
\begin{align*}
 P_Xu = \langle X,\nabla u\rangle + \frac{u}2\diver X
\end{align*}
satisfies
\begin{align*}
 u P_Xu = \diver\biggl(\frac12u^2X\biggr),
\end{align*}
we arrive at
\begin{align}
 \int_D (\Delta u-\lambda u) P_Xu
 = &\int_D \biggl(\langle\nabla_{\nabla u}X,\nabla u\rangle + \frac14\diver X\Delta u^2\biggr) \notag\\
 &- \int_{\partial D}\biggl(\langle X,\nabla u\rangle\langle\nabla u,\nu\rangle
 - \frac12\bigl(|\nabla u|^2 - \lambda u^2\bigr)\langle X,\nu\rangle\biggr),\label{x1}
\end{align}
which corresponds to~\cite[identity~(2)]{Xavier1988}, but with boundary integral included.
It will also be useful to have the latter formula with one of the substitutions
\begin{align}
 \int_D\diver X\Delta u^2
 &= \int_D\big\langle\nabla\diver X,\nabla u^2\big\rangle - \int_{\partial D}\diver X\nabla_\nu u^2 \label{x2a}\\
 &= \int_D(\Delta\diver X)u^2
 + \int_{\partial D} \bigl((\nabla_\nu\diver X)u^2 - \diver X\nabla_\nu u^2\bigr), \nonumber%\label{x2b}
\end{align}
which is Green's formula applied to the functions $\diver X$ and $u^2$.
Clearly,
\begin{gather}
 \int_D (\Delta u-\lambda u) P_Xu
	\le \|\Delta u-\lambda u\|_2\|P_Xu\|_2
	\le C\|\Delta u-\lambda u\|_2(\|\Delta u\|_2+\|u\|_2)\label{x3}
\end{gather}
if $X$ and $\diver X$ are bounded on $D$,
and $u$ satisfies a non-positive boundary condition.
This will be important in view of~\eqref{Kato9}.
